\hypertarget{relation-to-other-work}{%
\section{Relation to other work}\label{relation-to-other-work}}

\textbf{Unimodular gravity and vacuum-energy sequestering.} The
allocation model of Section 4 reproduces, in a finite setting, the
pattern in which a global constraint with a Lagrange multiplier keeps a
uniform energy offset from affecting the state
\citep{HenneauxTeitelboim1989, KaloperPadilla2014}. The closure-price
theorem adds a general condition for when such a multiplier is positive,
and the gauge theorem of Section 6.4 shows that in a closed-ledger model
the absolute value of the multiplier is not recoverable from the
invariant record. Like unimodular gravity, neither result fixes the
value.

\textbf{Anthropic and measure arguments.} The record window has the
shape of Weinberg's bound \citep{Weinberg1987}: structure formation
limits \(\lambda\) from above without placing the observed value near
the limit. The measure probe makes the dependence on the prior explicit
\citep{MartelShapiroWeinberg1998, BoussoPolchinski2000}. Section 7.3
proves that smoothness near zero does not guarantee concentration, and
that reweighting only the frozen base-grid values, with full support and
the stipulated conditional seed weights kept, cannot meet the fixed
top-decade concentration criterion. The window and its first-record
ratio bounds are exact theorems about the model; the prior remains an
input.

\textbf{Backreaction and averaging.} A companion pre-registration
\citep{TsiokosDarkEnergy2026} treats dark energy as a correction that
appears when a homogeneous completion fails to commute with
inhomogeneous evolution, and the audits of Section 8 are built from its
discriminators. The growth result is relevant to any program looking for
dynamical signatures in structure data: on the surveyed probes the
decisive within-data-set comparison is not available, and an apparent
strong-regime residual is sensitive to nonlinear modeling.

\textbf{Holography.} Section 4.4 states the precise finite relation
between min-cut area and an optimization price. Both are outputs of a
duality, but area is an optimal value under a declared scaling and the
closure price is a multiplier; the proposed shared composition rule for
entanglement and area fails on an explicit example. We use the relation
only as an analogy.

\textbf{The Six Birds program.} Earlier work applied the same calculus
to the structure of the Standard Model and to the relation between
quantum mechanics and general relativity \citep{TsiokosSixBirds2026}.
That work expected questions about contingent quantities to yield
windows and statements about what is undetermined, rather than value
laws, and questions about budgets to yield structural relations. The
results here fit that expectation: the budget question gave a theorem
about prices, and the value question gave readout-relative
nonidentifiability and a gauge result. They do not prove it, since the
value question was answered for specific readouts and a specific ledger,
not for all possible laws. This study also reused the earlier framework
deliberately, so it tests the framework's consistency and reach rather
than providing an independent confirmation.
